% =============================================================================
% Q4 -- detailed design requirements and technical specifications
% =============================================================================
\clearpage
\sectiontitle{Q4. Detailed Design Requirements and Technical Specifications}

\hypertarget{q4-requirement-basis}{}
\qtwosubsection{4.1 Requirement basis and evidence boundary}

{\fontsize{9.5}{11.2}\selectfont
Q4 converts the functional blocks isolated in Q3 into testable engineering
requirements. A specification is included only when its evidence status is
clear: \textbf{P} denotes a result demonstrated in a published precedent,
\textbf{D} denotes a value selected for the present design, and \textbf{V}
denotes a value that must be verified on the fabricated slab. A paper can
support a mechanism or a starting process without proving the performance of
the present ink--printer--PDMS--sensor combination.\par}

\evidencebox{Critical evidence boundary:}{No reviewed paper reports the exact
56\,$\times$\,55\,$\times$\,5.2~mm slab with directly printed silver
pressure/shear, temperature, and humidity sensors. The 5.2~mm thickness is a
project CAD and molding requirement. Published work separately supports molded
PDMS, direct Ag printing, multimodal prosthetic sensing, and prosthetic-liner
thickness effects; Q5 must verify their integration.}

\vspace{0.45em}
\begingroup
\fontsize{7.8}{9.15}\selectfont
\setlength{\tabcolsep}{3pt}\renewcommand{\arraystretch}{1.06}
\begin{tabularx}{\linewidth}{>{\raggedright\arraybackslash}p{3.35cm}Y >{\raggedright\arraybackslash\bfseries\color{AUBRed}}p{4.55cm}}
\rowcolor{AUBRed}
\color{white}\textbf{Q1 origin $\rightarrow$ linked Q2 evidence} &
\color{white}\textbf{Requirement purpose} &
\color{white}\textbf{Q4 specification families} \\
\textbf{Q1-1} $\rightarrow$ \hyperlink{q2-functional-blocks}{\textcolor{AUBRed}{\bfseries\underline{Q2 \S2.1}}} / \hyperlink{q2-synthesis}{\textcolor{AUBRed}{\bfseries\underline{\S2.8}}} &
Four-mode integration in the 5.2~mm slab & GEO, MAT, PRT, SYS \\
\rowcolor{AUBPale}
\textbf{Q1-2} $\rightarrow$ \hyperlink{q2-humidity}{\textcolor{AUBRed}{\bfseries\underline{Q2 \S2.4.1}}} &
Humidity vapour access and liquid protection & HUM, ENV, SAFE \\
\textbf{Q1-3} $\rightarrow$ \hyperlink{q2-pressure-shear}{\textcolor{AUBRed}{\bfseries\underline{Q2 \S2.3}}} / \hyperlink{q2-temperature}{\textcolor{AUBRed}{\bfseries\underline{\S2.4}}} / \hyperlink{q2-humidity}{\textcolor{AUBRed}{\bfseries\underline{\S2.4.1}}} &
Detection of controlled interface conditions & MECH, TEMP, HUM \\
\rowcolor{AUBPale}
\textbf{Q1-4} $\rightarrow$ \hyperlink{q2-current-build}{\textcolor{AUBRed}{\bfseries\underline{Q2 \S2.2}}} / \hyperlink{q2-pdms-printing}{\textcolor{AUBRed}{\bfseries\underline{\S2.5}}} &
Manufacturability, available equipment, and cost & MFG \\
\textbf{Q1-5} $\rightarrow$ \hyperlink{q2-whole-liner}{\textcolor{AUBRed}{\bfseries\underline{Q2 \S2.7}}} / \hyperlink{q2-synthesis}{\textcolor{AUBRed}{\bfseries\underline{\S2.8}}} &
Synchronized coordinate mapping and interpretation & DAQ, MAP, ALG \\
\rowcolor{AUBPale}
\textbf{Q1-6} $\rightarrow$ \hyperlink{q2-current-build}{\textcolor{AUBRed}{\bfseries\underline{Q2 \S2.2}}} / \hyperlink{q2-whole-liner}{\textcolor{AUBRed}{\bfseries\underline{\S2.7}}} &
Sensor density, coverage, and localization & MAP \\
\textbf{Q1-7} $\rightarrow$ \hyperlink{q2-pdms-printing}{\textcolor{AUBRed}{\bfseries\underline{Q2 \S2.5}}} / \hyperlink{q2-hydrogel-interactions}{\textcolor{AUBRed}{\bfseries\underline{\S2.6.1}}} &
Drift, environmental exposure, and durability & DUR, ENV \\
\end{tabularx}
\endgroup
{\fontsize{7.55}{8.7}\selectfont\color{AUBGray}
\textbf{Table Q4-1. Clickable Q1--Q2--Q4 traceability.} Q1 identifies the
requirement origin; each underlined red Q2 reference opens the detailed design,
literature, or implementation evidence used to define its Q4 specification.\par}

\vspace{0.55em}
\begingroup
\fontsize{8.4}{9.85}\selectfont
\setlength{\tabcolsep}{3pt}
\renewcommand{\arraystretch}{1.12}
\begin{tabularx}{\linewidth}{>{\raggedright\arraybackslash\bfseries\color{AUBRed}}p{1.15cm}Y Y Y}
\rowcolor{AUBRed}
\color{white}\textbf{ID} & \color{white}\textbf{Project requirement} &
\color{white}\textbf{Basis and status} & \color{white}\textbf{Acceptance evidence} \\
GEO-01 & Active slab envelope shall be 56\,$\times$\,55~mm with nominal total
thickness 5.2~mm. & \textbf{D:} fixed by current CAD/mold. Published 4 and
6~mm liner models provide context only \cite{q4addacherif2025}. & Dimensioned
CAD export; caliper thickness at centre and four corners. \\
\rowcolor{AUBPale}
GEO-02 & Total thickness shall be $5.2\pm0.2$~mm; five-point range shall not
exceed 0.4~mm. & \textbf{D/V:} preliminary manufacturing tolerance, not a
literature-derived clinical limit. & Five recorded measurements for every
slab; reject visible voids, incomplete fill, or out-of-tolerance specimens. \\
GEO-03 & Baseline stack shall allocate 0.20~mm skin-side cover, 3.50~mm sensor
carrier, and 1.50~mm back encapsulation, summing to 5.20~mm. & \textbf{D/V:}
current engineering allocation; layer transfer and local stiffness require
verification. & Cross-sectional coupon or mold witness; report each layer and
sensor-plane depth rather than total thickness alone. \\
\rowcolor{AUBPale}
MAT-01 & The first calibrated build shall use bare PDMS as H0; hydrogel is not
part of the baseline specification. & \textbf{D:}
\hyperlink{q2-hydrogel-study}{\textcolor{AUBRed}{\underline{Q2 Section 2.6 evidence gate}}}. Hydrogel
changes moisture transport and stress transfer \cite{yuk2016,lu2025}. & H0
calibration completed before any H1/H2 material comparison. \\
\end{tabularx}
\endgroup
\sourceanchor{Thickness context from Adda Cherif \textit{et al.}
\cite{q4addacherif2025}; dimensions and tolerances are explicitly identified
as project decisions rather than published outcomes.}

\clearpage
\hypertarget{q4-implementation-geometry}{}
\sectiontitle{Q4. Implementation Geometry and Liner Context}

\qtwosubsection{4.2 Project slab, sensor placement, and thickness precedent}

{\fontsize{9.25}{10.95}\selectfont
The two supplied CAD views are implementation inputs: the render communicates
the physical arrangement, while the blueprint supplies the coordinate and
enclosure reference used by MAP-01. They are shown here with the published
liner-thickness precedent because Q4 defines how the geometry becomes a
measurable specification. Neither perspective view replaces a dimensioned
manufacturing drawing.\par}

\begin{minipage}[t]{0.485\textwidth}
  \vspace{0pt}\centering
  \includegraphics[width=\linewidth,height=63mm,keepaspectratio]
    {assets/q2/q2-a_slab_design_render.png}\par
  {\fontsize{7.6}{8.75}\selectfont\color{AUBGray}\raggedright
  \textbf{Figure Q4-1A. Slab design render.} Supplied three-dimensional view
  of the distributed sensing elements and central enclosure region.\par}
\end{minipage}\hfill
\begin{minipage}[t]{0.485\textwidth}
  \vspace{0pt}\centering
  \includegraphics[width=\linewidth,height=63mm,keepaspectratio]
    {assets/q2/q2-d_sensor_placement_blueprint.png}\par
  {\fontsize{7.6}{8.75}\selectfont\color{AUBGray}\raggedright
  \textbf{Figure Q4-1B. Sensor-placement blueprint.} Supplied CAD view showing
  relative sensor placement, the enclosure principle, and conductor exits.\par}
\end{minipage}

\vspace{0.55em}
\begin{center}
  \includegraphics[width=0.72\linewidth,height=61mm,keepaspectratio]
    {paper_sources/figures_selected/fem_fig2_liner_thickness_models.png}\par
\end{center}
{\fontsize{7.7}{8.9}\selectfont\color{AUBGray}
\textbf{Figure Q4-1C. Published liner-thickness context.} Adda Cherif
\textit{et al.} compared 2, 4, and 6~mm transfemoral liner models
\cite{q4addacherif2025}. The study brackets the project's 5.2~mm value but does
not validate the printed sensor stack or the exact selected thickness. CC BY
4.0.\par}

\vspace{0.45em}
\evidencebox{Required drawing output:}{Release to fabrication requires one
orthographic plan and one cross-section labelled with the 56\,$\times$\,55~mm
envelope, 5.2~mm total thickness, layer depths, sensor IDs and footprints,
coordinate origin, conductor widths, pad positions, edge clearances, enclosure
boundary, and nominal versus measured dimensions.}

\clearpage
\sectiontitle{Q4. PDMS Molding and Direct Silver Printing}

\hypertarget{q4-fabrication}{}
\qtwosubsection{4.3 A literature-backed fabrication window}

{\fontsize{9.15}{10.8}\selectfont
The closest direct precedent is not a 5.2~mm prosthetic slab. Abu-Khalaf
\textit{et al.} molded 1~mm Sylgard~184 specimens in acrylic molds, using a
10:1 base-to-curing-agent ratio, 30~min vacuum degassing, and 70~$^{\circ}$C
curing for 2~h before direct AgNP printing \cite{q4abukhalaf2018}. Albrecht
\textit{et al.} established a separate low-temperature printing route using
oxygen-plasma-treated PDMS, printing within 30~min, and drying at
60~$^{\circ}$C for 30~min \cite{albrecht2018stretch}. These are \emph{separate
published starting recipes}; their parameters must not be combined silently.\par}

\vspace{0.25em}
\begin{center}
\begin{tikzpicture}[x=1cm,y=1cm,>=stealth]
  \tikzstyle{source}=[draw=AUBLine,fill=AUBPale,rounded corners=1pt,
    text width=44mm,minimum height=9mm,align=left,
    font=\fontsize{6.5}{7.4}\selectfont]
  \tikzstyle{action}=[draw=AUBRed!60,fill=white,rounded corners=1pt,
    text width=35mm,minimum height=10mm,align=center,
    font=\fontsize{6.7}{7.6}\selectfont]
  \node[source] (a) at (2.4,3.2) {\textbf{Molding precedent}\quad 10:1
    Sylgard 184; 30~min vacuum; 70~$^{\circ}$C, 2~h; 1~mm coupon
    \cite{q4abukhalaf2018}};
  \node[source] (b) at (2.4,1.8) {\textbf{Surface/print precedent}\quad oxygen
    plasma; print within 30~min; 60~$^{\circ}$C, 30~min
    \cite{albrecht2018stretch}};
  \node[source] (c) at (2.4,0.4) {\textbf{Geometry precedent}\quad
    $\sim$1.8~mm interconnect and horseshoe DOE; precure-dependent embedding
    \cite{q4abukhalaf2019,sun2016}};
  \node[action] (doe) at (7.6,1.8) {\textbf{Project coupon matrix}\\material
    recipe $\times$ treatment $\times$ print delay $\times$ trace geometry};
  \node[action] (select) at (12.1,1.8) {\textbf{Select one recorded window}\\
    yield, dimensions, resistance, adhesion, insulation};
  \node[action] (slab) at (16.5,1.8) {\textbf{Release to 5.2~mm slab}\\only
    after coupon acceptance};
  \draw[->,thick,color=AUBRed] (a.east)--(doe.west);
  \draw[->,thick,color=AUBRed] (b.east)--(doe.west);
  \draw[->,thick,color=AUBRed] (c.east)--(doe.west);
  \draw[->,thick,color=AUBRed] (doe)--(select);
  \draw[->,thick,color=AUBRed] (select)--(slab);
\end{tikzpicture}
\end{center}
{\fontsize{7.4}{8.55}\selectfont\color{AUBGray}
\textbf{Figure Q4-2. Evidence-to-process decision route.} The published routes
remain separate inputs. A project coupon matrix selects one compatible process
before the full slab is printed; this is an implementation diagram, not a
reproduction of the paper figures already shown in Q2.\par}

\vspace{0.3em}
\begingroup
\fontsize{8.0}{9.35}\selectfont
\setlength{\tabcolsep}{3pt}\renewcommand{\arraystretch}{1.07}
\begin{tabularx}{\linewidth}{>{\raggedright\arraybackslash\bfseries\color{AUBRed}}p{1.2cm}Y Y Y}
\rowcolor{AUBRed}
\color{white}\textbf{ID} & \color{white}\textbf{Published precedent} &
\color{white}\textbf{Project specification} & \color{white}\textbf{Pass evidence} \\
MFG-01 & Sylgard~184, 10:1, 30~min vacuum, 70~$^{\circ}$C for 2~h on 1~mm
specimens \cite{q4abukhalaf2018}. & Use as the documented candidate recipe if
the current PDMS is Sylgard~184; otherwise record material, ratio, degassing,
cure, and hardness as a separate recipe. & Complete fill; no visible bubbles in
active areas; GEO-01/02 measurements pass. \\
\rowcolor{AUBPale}
PRT-01 & Continuous AgNP lines required plasma; printing occurred within
30~min and drying was 60~$^{\circ}$C for 30~min
\cite{albrecht2018stretch}. & Run untreated and treated coupons with fixed
treatment-to-print delay before committing the slab. & Printed width/profile,
open/short count, sheet resistance, and yield reported by treatment. \\
PRT-02 & Reliable low-cost-printer interconnects were 1--2~mm; an independent
DOE found a 1.78~mm straight trace optimum and 4~mm-amplitude horseshoe traces
for strain tolerance \cite{albrecht2018stretch,q4abukhalaf2019}. & Retain the
functional sensor fingers from their source designs; use $\geq1.8$~mm where
space permits for long interconnects. & 100\% intended paths continuous, no
adjacent shorts, and printed dimensions within $\pm10\%$ of artwork. \\
\rowcolor{AUBPale}
PRT-03 & Bare nanostructured Ag traces failed tape adhesion; thin coating or a
PDMS sandwich was recommended \cite{albrecht2018stretch}. & No skin-contacting
silver. Encapsulate all traces except the controlled RH vapour window. & Visual
inspection, dry/wet insulation test, adhesion test, and pre/post resistance. \\
\end{tabularx}
\endgroup
\sourceanchor{Direct molding and Ag-printing parameters are traceable to
Abu-Khalaf \textit{et al.} and Albrecht \textit{et al.}
\cite{q4abukhalaf2018,q4abukhalaf2019,albrecht2018stretch}; the full 5.2~mm
process remains a project verification item.}

% -----------------------------------------------------------------------------
% Supplied PDMS mixing tutorial
% Source: PDMS_Mixing_Tutorial.docx
% Wording and source-image order are preserved.
% -----------------------------------------------------------------------------

\clearpage
\hypertarget{q4-pdms-mixing-tutorial}{}
\sectiontitle{Q4. PDMS Mixing Tutorial (Polydimethylsiloxane)}
\qtwosubsection{Practical mixing procedure}

\begingroup
\fontsize{9.0}{10.45}\selectfont
\newcommand{\pdmshead}[1]{\par\vspace{0.48em}{\bfseries #1}\par\vspace{0.15em}}

This tutorial describes the general process for mixing PDMS silicone elastomer
for molding, coatings, and laboratory applications.

\pdmshead{Materials}
\begin{itemize}[leftmargin=1.45em,itemsep=0.15em,topsep=0.25em]
  \item PDMS base (silicone prepolymer)
  \item Curing agent (crosslinker)
  \item Mixing cup and spatula
  \item Balance or measuring tools
  \item Vacuum desiccator (for degassing)
  \item Oven or hot plate for curing
\end{itemize}

\pdmshead{Step 1: Measure the Components}
We used Sylgard 184 mixing ratio which is 10:1 by weight (base:curing agent).
Example: 10 g PDMS base mixed with 1 g curing agent. Different ratios can change
the softness, curing time, and mechanical properties.

\pdmshead{Step 2: Mix Thoroughly}
Add the curing agent to the PDMS base and mix slowly for 5--10 minutes. Scrape
the sides and bottom of the container during mixing. Avoid vigorous stirring
because it introduces air bubbles.

\pdmshead{Step 3: Degas the PDMS}
Place the mixed PDMS in a vacuum chamber. Apply vacuum until trapped air bubbles
expand and disappear. This usually takes 15--60 minutes depending on the volume
and viscosity. Release the vacuum slowly to prevent overflow.

\pdmshead{Step 4: Pour or Coat}
Transfer the degassed PDMS onto the mold or substrate carefully. Pour slowly to
reduce the chance of trapping new air bubbles.

\pdmshead{Step 5: Cure the PDMS}
Typical Sylgard 184 curing conditions are 60--80~$^{\circ}$C for 1--2 hours.
PDMS can also cure at room temperature, but this may require about 24 hours or
longer depending on conditions.

\pdmshead{Step 6: Demold and Finish}
Allow the PDMS to cool after curing. Carefully remove it from the mold and trim
edges or excess material as required.

\pdmshead{Important Tips}
\begin{itemize}[leftmargin=1.45em,itemsep=0.15em,topsep=0.25em]
  \item More curing agent generally produces a harder and faster-curing PDMS.
  \item Less curing agent generally produces softer and more flexible PDMS.
  \item Proper degassing is important for removing bubbles.
  \item Use clean tools because contamination may affect curing.
\end{itemize}
\endgroup

\clearpage
\sectiontitle{Q4. PDMS Mixing Tutorial -- Laboratory Images}
\begin{center}
  \centering
  \includegraphics[angle=270,origin=c,width=0.47\linewidth,height=0.36\textheight,keepaspectratio]{assets/pdms_mixing/01_balance_setup.jpeg}\hfill
  \includegraphics[angle=270,origin=c,width=0.47\linewidth,height=0.36\textheight,keepaspectratio]{assets/pdms_mixing/02_component_handling.jpeg}

  \vspace{0.8em}
  \includegraphics[angle=270,origin=c,width=0.47\linewidth,height=0.36\textheight,keepaspectratio]{assets/pdms_mixing/03_curing_agent_mass.jpeg}\hfill
  \includegraphics[angle=270,origin=c,width=0.47\linewidth,height=0.36\textheight,keepaspectratio]{assets/pdms_mixing/04_pdms_base_mass.jpeg}
\end{center}

\clearpage
\sectiontitle{Q4. PDMS Mixing Tutorial -- Laboratory Images}
\begin{center}
  \centering
  \includegraphics[angle=270,origin=c,width=0.47\linewidth,height=0.70\textheight,keepaspectratio]{assets/pdms_mixing/05_manual_mixing.jpeg}\hfill
  \includegraphics[angle=270,origin=c,width=0.47\linewidth,height=0.70\textheight,keepaspectratio]{assets/pdms_mixing/06_vacuum_oven.jpeg}
\end{center}


\clearpage
\sectiontitle{Q4. Pressure and Directional-Shear Requirements}

\hypertarget{q4-mechanical}{}
\qtwosubsection{4.4 Mechanical sensing envelope, accuracy, and separation}

{\fontsize{9.2}{10.85}\selectfont
Albrecht \textit{et al.} demonstrated the selected four-capacitor topology on
170~$\mu$m PET with a 40~$\mu$m PDMS dielectric---not on a 5.2~mm molded PDMS
slab. The reported cell separated normal and two shear directions over
0.1--8~N, with approximately 5.2~fF/N normal and 13.1~fF/N shear sensitivity,
but also showed viscoelastic drift and shear hysteresis
\cite{albrecht2019}. The topology is therefore reusable; its calibration is
not.\par}

\begin{minipage}[t]{0.61\textwidth}
  \vspace{0pt}\centering
  \begin{tikzpicture}[x=0.86cm,y=0.016cm]
    \fill[AUBGreen!10] (0,0) rectangle (10,100);
    \fill[AUBYellow!10] (0,100) rectangle (10,250);
    \draw[->,thick] (0,0)--(10.7,0);
    \draw[->,thick] (0,0)--(0,270);
    \node[below=3.8mm,font=\fontsize{6.7}{7.5}\selectfont]
      at (5,0) {signed shear (kPa)};
    \node[rotate=90,font=\fontsize{6.7}{7.5}\selectfont]
      at (-0.85,125) {normal pressure (kPa)};
    \foreach \x/\lab in {0/-50,2.5/-25,5/0,7.5/25,10/50}{
      \draw[AUBLine] (\x,0)--(\x,250);
      \node[below,font=\fontsize{6.3}{7}\selectfont] at (\x,0) {\lab};
    }
    \foreach \y in {0,25,50,75,100,150,200,250}{
      \draw[AUBLine] (0,\y)--(10,\y);
      \node[left,font=\fontsize{6.3}{7}\selectfont] at (0,\y) {\y};
    }
    \foreach \x in {0,2.5,5,7.5,10}{
      \foreach \y in {0,25,50,75,100,150,200,250}{
        \fill[AUBRed] (\x,\y) circle (1.6pt);
      }
    }
    \node[anchor=north east,font=\fontsize{6.4}{7.2}\selectfont\bfseries,
      text=AUBGreen] at (9.8,96) {routine band};
    \node[anchor=north east,font=\fontsize{6.4}{7.2}\selectfont\bfseries,
      text=AUBYellow!70!black] at (9.8,246) {extended characterization};
  \end{tikzpicture}
\end{minipage}\hfill
\begin{minipage}[t]{0.35\textwidth}
  \vspace{3mm}
  \begin{tikzpicture}[>=stealth,node distance=2.4mm]
    \tikzstyle{mechstep}=[draw=AUBLine,fill=AUBPale,rounded corners=1pt,
      text width=42mm,minimum height=7mm,align=center,
      font=\fontsize{6.7}{7.6}\selectfont]
    \node[mechstep] (s1) {set normal preload};
    \node[mechstep,below=of s1] (s2) {apply signed $x/y$ shear};
    \node[mechstep,below=of s2] (s3) {dwell 5, 30, 300~s};
    \node[mechstep,below=of s3] (s4) {unload and record recovery};
    \node[mechstep,below=of s4] (s5) {reverse direction; repeat};
    \draw[->,color=AUBRed,thick] (s1)--(s2);
    \draw[->,color=AUBRed,thick] (s2)--(s3);
    \draw[->,color=AUBRed,thick] (s3)--(s4);
    \draw[->,color=AUBRed,thick] (s4)--(s5);
  \end{tikzpicture}
\end{minipage}
{\fontsize{7.4}{8.55}\selectfont\color{AUBGray}
\textbf{Figure Q4-3. Planned mechanical characterization matrix.} Red points
are required pressure--shear combinations, not measured results. The envelope
is selected from the interface ranges reviewed in Q3
\cite{q3sanders1992,q3tang2022,q3zhang2001}; the sequence exposes hysteresis,
cross-axis response, and PDMS recovery.\par}

\vspace{0.5em}
\begingroup
\fontsize{8.2}{9.6}\selectfont
\setlength{\tabcolsep}{3pt}\renewcommand{\arraystretch}{1.11}
\begin{tabularx}{\linewidth}{>{\raggedright\arraybackslash\bfseries\color{AUBRed}}p{1.25cm}Y Y Y}
\rowcolor{AUBRed}
\color{white}\textbf{ID} & \color{white}\textbf{Requirement} &
\color{white}\textbf{Why this value was selected} & \color{white}\textbf{Verification and gate} \\
MECH-01 & Characterize normal pressure from 0 to 250~kPa; designate 0--100~kPa
as the routine operating band. & \textbf{D:} published interfaces include
measured tens-of-kPa values and a reported 226~kPa finite-element peak
\cite{q3sanders1992,q3tang2022,q3zhang2001}. These are context, not universal
injury limits. & Dead weights/load cell and known footprint; report full-range
and routine-band calibration, residuals, saturation, and overload recovery. \\
\rowcolor{AUBPale}
MECH-02 & Characterize signed shear from $-50$ to $+50$~kPa in both slab axes
under at least three normal preloads. & \textbf{D:} selected to include
published shear values up to roughly 40--50~kPa
\cite{q3sanders1992,q3zhang2001}. & Independent tangential force/displacement
reference; report vector magnitude, angle, cross-axis coupling, and slip onset. \\
MECH-03 & Within each calibrated band: $R^2\geq0.95$, NRMSE $\leq10\%$,
hysteresis $\leq10\%$ full scale, repeatability CV $\leq10\%$. & \textbf{D/V:}
predeclared engineering gates; not clinical limits and not inherited from the
published PET cell. & At least three independently fabricated slabs, increasing
and decreasing steps, three cycles per level, and held-out verification points. \\
\rowcolor{AUBPale}
MECH-04 & Non-target axis response shall be $\leq15\%$ of target full-scale
response; baseline shall recover to within 10\% full scale after unloading. &
\textbf{D/V:} matches Q3 separation and recovery gates. PDMS viscoelasticity
makes both checks necessary \cite{albrecht2019}. & Modality-by-modality matrix;
5, 30, and 300~s dwell/recovery records; failure cannot be averaged into a
combined score. \\
MECH-05 & Shear signals may indicate a possible slip event but shall not label
slip without an independent relative-motion reference. & \textbf{P/D:} force
and pistoning are related but distinct measurements \cite{q3tang2022,noll2018}. &
Synchronized camera/encoder trace; onset-time error and confusion matrix on
held-out repetitions. \\
\end{tabularx}
\endgroup
\sourceanchor{Mechanical envelope drawn from prosthetic-interface literature;
sensor topology and its transfer limitation from Albrecht \textit{et al.}
\cite{albrecht2019,q3sanders1992,q3tang2022,q3zhang2001}.}

\clearpage
\sectiontitle{Q4. Temperature and Humidity Requirements}

\hypertarget{q4-temp-humidity}{}
\qtwosubsection{4.5 Separate measurands, separate protection, common coordinates}

{\fontsize{8.75}{10.25}\selectfont
The source temperature meander and humidity IDE were demonstrated on PET and
Kapton, respectively---not on molded PDMS. Their geometry and response ranges
are starting specifications; transfer, protection, and calibration remain
project verification items \cite{liew2020,aitmammar2020}.\par}

\begin{minipage}[t]{0.485\textwidth}
  \vspace{0pt}\centering
  \begin{tikzpicture}
  \begin{axis}[
    width=\linewidth,height=38mm,
    xmin=0,xmax=10,ymin=20,ymax=50,
    xlabel={step index},ylabel={temperature ($^{\circ}$C)},
    xtick={0,2,4,6,8,10},ytick={25,30,35,40,45},
    grid=both,axis line style={AUBGray},tick label style={font=\fontsize{6}{6.8}\selectfont},
    label style={font=\fontsize{6.5}{7.3}\selectfont}]
    \addplot+[const plot,very thick,color=AUBRed,mark=*]
      coordinates {(0,25)(1,25)(1,30)(2,30)(2,35)(3,35)(3,40)(4,40)
      (4,45)(5,45)(5,40)(6,40)(6,35)(7,35)(7,30)(8,30)(8,25)(10,25)};
  \end{axis}
  \end{tikzpicture}
\end{minipage}\hfill
\begin{minipage}[t]{0.485\textwidth}
  \vspace{0pt}\centering
  \begin{tikzpicture}
  \begin{axis}[
    width=\linewidth,height=38mm,
    xmin=0,xmax=10,ymin=10,ymax=100,
    xlabel={step index},ylabel={relative humidity (\%RH)},
    xtick={0,2,4,6,8,10},ytick={20,40,60,80,95},
    grid=both,axis line style={AUBGray},tick label style={font=\fontsize{6}{6.8}\selectfont},
    label style={font=\fontsize{6.5}{7.3}\selectfont}]
    \addplot+[const plot,very thick,color=AUBBlue,mark=*]
      coordinates {(0,20)(1,20)(1,40)(2,40)(2,60)(3,60)(3,80)(4,80)
      (4,95)(5,95)(5,80)(6,80)(6,60)(7,60)(7,40)(8,40)(8,20)(10,20)};
  \end{axis}
  \end{tikzpicture}
\end{minipage}
{\fontsize{7.35}{8.45}\selectfont\color{AUBGray}
\textbf{Figure Q4-4. Planned chamber profiles.} Left: increasing, decreasing,
and recovery temperature levels. Right: the corresponding RH sequence, repeated
at two or more temperatures. These are prescribed reference inputs, not project
measurements; the levels follow the calibration ranges reviewed in Q2
\cite{liew2020,aitmammar2020,paterno2020}.\par}

\vspace{0.3em}
\begingroup
\fontsize{7.7}{9.0}\selectfont
\setlength{\tabcolsep}{3pt}\renewcommand{\arraystretch}{1.02}
\begin{tabularx}{\linewidth}{>{\raggedright\arraybackslash\bfseries\color{AUBRed}}p{1.2cm}Y Y Y}
\rowcolor{AUBRed}
\color{white}\textbf{ID} & \color{white}\textbf{Project specification} &
\color{white}\textbf{Published anchor} & \color{white}\textbf{Verification} \\
TEMP-01 & Calibrate every temperature channel from 25 to 45~$^{\circ}$C at
5~$^{\circ}$C steps; target error $\leq0.5$~$^{\circ}$C and displayed
resolution $\leq0.1$~$^{\circ}$C. & Liew used 30--100~$^{\circ}$C; Patern\`o
calibrated liner sensors over 25--45~$^{\circ}$C and demonstrated substantially
smaller post-correction error using different devices
\cite{liew2020,paterno2020}. & Co-located traceable RTD/thermocouple; increasing,
decreasing, and held-out temperatures; RMSE, bias, $t_{90}$, and recovery. \\
\rowcolor{AUBPale}
TEMP-02 & Use the project meander geometry (6.9\,$\times$\,13.2~mm, 0.5~mm
track, 0.4~mm gap) and quantify strain cross-sensitivity. & Geometry taken from
Liew \textit{et al.}; their sensitivities were 0.0543 and 0.1086~$\Omega$/
$^{\circ}$C for two inks \cite{liew2020}. & Compare resistance under temperature
only, strain only, and combined loading; published sensitivity is not the
project calibration coefficient. \\
HUM-01 & Calibrate vapour response from 20\% to 95\%RH at two or more
temperatures; target error $\leq5$\%RH, hysteresis $\leq10$\% full scale, and
vapour $t_{90}\leq100$~s before final protection selection. & Published CAB IDE
covered 20--90\%RH and reported response below 100~s and recovery below 5~s;
prosthetic liner calibration extended to 95\%RH with different devices
\cite{aitmammar2020,paterno2020}. & Reference hygrometer; increasing/decreasing
RH, held-out levels, temperature compensation, and dry recovery. \\
\rowcolor{AUBPale}
HUM-02 & The humidity region shall receive vapour through a hydrophobic,
vapour-permeable barrier; liquid shall not directly contact Ag traces or
electrical joints. & Patern\`o protected non-waterproof sensors with breathable
waterproof fabric; the printed IDE paper did not test prosthetic sweat
\cite{paterno2020,aitmammar2020}. & Vapour steps followed by controlled
artificial-sweat droplets; inspect leakage, saturation, baseline recovery, and
post-exposure calibration. \\
HUM-03 & Vapour RH and liquid sweat shall remain separate software states; a
saturated or nonrecovering IDE shall be flagged, not converted to a higher
``discomfort'' score. & \textbf{D/V:} measurement-quality requirement arising
from the different transport mechanisms. & Quality flag triggered by out-of-
range capacitance, recovery failure, leakage, or reference disagreement. \\
\end{tabularx}
\endgroup
\sourceanchor{Temperature and RH geometry/performance precedents from Liew,
A\"{\i}t-Mammar, and Patern\`o \textit{et al.}
\cite{liew2020,aitmammar2020,paterno2020}.}

\clearpage
\sectiontitle{Q4. Acquisition, Coordinates, and Interface-Condition Mapping}

\hypertarget{q4-mapping}{}
\qtwosubsection{4.6 Data integrity is part of the sensor specification}

{\fontsize{9.25}{10.9}\selectfont
The physical slab does not produce a meaningful map by itself. Every channel
must retain its units, calibration version, timestamp, and CAD coordinate
before interpolation. Kwak \textit{et al.} demonstrated a wireless prosthetic
interface at 30~Hz; Patern\`o \textit{et al.} logged slower temperature/RH
channels every 5~s; Turner \textit{et al.} compared interpolation methods on a
three-dimensional prosthetic socket \cite{q3kwak2020,paterno2020,q3turner2021}.
These precedents justify a multirate pipeline rather than forcing every
modality to the same physical response speed.\par}

\begin{center}
\begin{tikzpicture}[>=stealth,node distance=3mm]
  \tikzstyle{data}=[draw=AUBLine,fill=AUBPale,rounded corners=1pt,
    minimum height=9mm,text width=31mm,align=center,
    font=\fontsize{6.9}{7.9}\selectfont]
  \node[data] (d1) {raw channels\\and reference};
  \node[data,right=of d1] (d2) {calibration +\\quality flags};
  \node[data,right=of d2] (d3) {sensor ID $\rightarrow$\\CAD coordinate};
  \node[data,right=of d3] (d4) {separate physical-\\unit maps};
  \node[data,right=of d4] (d5) {optional sourced\\condition score};
  \draw[->,thick,color=AUBRed] (d1)--(d2);
  \draw[->,thick,color=AUBRed] (d2)--(d3);
  \draw[->,thick,color=AUBRed] (d3)--(d4);
  \draw[->,thick,color=AUBRed] (d4)--(d5);
  \draw[->,dashed,color=AUBGray] (d2.south) to[bend right=25]
    node[below,font=\fontsize{6.2}{7}\selectfont]{failed quality gate stops interpretation} (d5.south);
\end{tikzpicture}
\end{center}
{\fontsize{7.5}{8.7}\selectfont\color{AUBGray}
\textbf{Figure Q4-5. Traceable mapping pipeline.} A smooth image is the last
step, not the evidence source; failed calibration or coordinate registration
cannot be hidden by interpolation.\par}

\vspace{0.45em}
\begingroup
\fontsize{8.2}{9.6}\selectfont
\setlength{\tabcolsep}{3pt}\renewcommand{\arraystretch}{1.11}
\begin{tabularx}{\linewidth}{>{\raggedright\arraybackslash\bfseries\color{AUBRed}}p{1.25cm}Y Y Y}
\rowcolor{AUBRed}
\color{white}\textbf{ID} & \color{white}\textbf{Requirement} &
\color{white}\textbf{Literature/design basis} & \color{white}\textbf{Acceptance evidence} \\
DAQ-01 & Pressure/shear shall be sampled at $\geq50$~Hz; temperature/RH at
$\geq0.2$~Hz. All records shall share one monotonic clock. & \textbf{D:}
50~Hz gives margin over a published 30~Hz prosthetic system; 0.2~Hz follows a
5~s liner microclimate precedent \cite{q3kwak2020,paterno2020}. & Timestamp
audit; channel rate, jitter, and lost samples reported during a 30~min run. \\
\rowcolor{AUBPale}
DAQ-02 & Inter-channel time skew shall be $\leq20$~ms for pressure/shear; packet
loss shall be $<1\%$; raw and filtered values shall both be retained. &
\textbf{D/V:} engineering synchronization and traceability targets. & Inject a
common electrical/mechanical event; verify timestamps, sample count, and filter
delay. \\
MAP-01 & Each channel shall have a unique ID, modality, unit, slab ID, CAD
$(x,y,z)$ coordinate, active footprint, orientation, and calibration version. &
Personalized liner work demonstrates the importance of registered placement
\cite{paterno2020}; this schema is project-specific. & Automated manifest-to-CAD
check: zero duplicate IDs, missing units, or unassigned coordinates. \\
\rowcolor{AUBPale}
MAP-02 & Known hotspot median localization error shall not exceed half the
nearest-neighbour sensor pitch; the 95th percentile and boundary errors shall
also be reported. & \textbf{D/V:} Q3 preregistered gate; Turner demonstrates
interpolation choices but not this array's accuracy \cite{q3turner2021}. &
Blinded stimulus positions; true/detected coordinates, error vectors, and
confidence/coverage mask. \\
MAP-03 & Display pressure, signed shear, temperature, and RH separately before
any fusion. Show sensor markers and measured coverage; do not extrapolate a
clinical conclusion outside validated coverage. & Prosthetic maps can be
smoothly interpolated, but interpolation is not measurement
\cite{q3turner2021}. & Screenshot plus numerical test: known inputs, correct
units, correct coordinates, no value outside coverage without a flag. \\
\rowcolor{AUBPale}
ALG-01 & Any combined interface-condition score shall publish its equation,
normalization range, threshold source, version, and missing-channel rule. It
shall not be labelled a validated discomfort-risk score without participant
evidence. & \textbf{D:} claim boundary established in Q1--Q3 and supported by
the variability reported in interface-pressure reviews
\cite{q3armitage2026}. & Software unit tests, versioned configuration, and a
display disclaimer until participant validation is completed. \\
\end{tabularx}
\endgroup
\sourceanchor{Sampling precedents from Kwak and Patern\`o; spatial mapping
precedent from Turner \textit{et al.}
\cite{q3kwak2020,paterno2020,q3turner2021}.}

\clearpage
\sectiontitle{Q4. Durability, Safety, Manufacturability, and Traceability}

\hypertarget{q4-design-gate}{}
\qtwosubsection{4.7 Master design gate before Q5 testing}

\begin{center}
\vspace{-0.2em}
\begin{tikzpicture}[>=stealth,node distance=3mm]
  \tikzstyle{gate}=[draw=AUBLine,fill=AUBPale,rounded corners=1pt,
    minimum height=7mm,text width=36mm,align=center,
    font=\fontsize{7}{8}\selectfont]
  \node[gate] (g1) {coupon process\\gate};
  \node[gate,right=of g1] (g2) {flat 5.2~mm slab\\calibration gate};
  \node[gate,right=of g2] (g3) {environmental +\\cyclic gate};
  \node[gate,right=of g3] (g4) {curved phantom /\\liner gate};
  \draw[->,thick,color=AUBRed] (g1)--(g2);
  \draw[->,thick,color=AUBRed] (g2)--(g3);
  \draw[->,thick,color=AUBRed] (g3)--(g4);
\end{tikzpicture}
\end{center}
{\fontsize{7.25}{8.4}\selectfont\color{AUBGray}
\textbf{Figure Q4-6. Build-and-test release sequence.} A failed process coupon
or flat-slab calibration blocks curved integration. Published benchmarks
include 1,000-cycle prosthetic-sensor and $>10{,}000$-cycle printed-PDMS
conductor tests \cite{q3kwak2020,albrecht2018stretch}.\par}

\vspace{0.25em}
\begingroup
\fontsize{8.0}{9.35}\selectfont
\setlength{\tabcolsep}{3pt}\renewcommand{\arraystretch}{1.06}
\begin{tabularx}{\linewidth}{>{\raggedright\arraybackslash\bfseries\color{AUBRed}}p{1.25cm}Y Y Y}
\rowcolor{AUBRed}
\color{white}\textbf{ID} & \color{white}\textbf{Requirement} &
\color{white}\textbf{Evidence basis} & \color{white}\textbf{Release criterion} \\
DUR-01 & Stage cyclic characterization at 1,000 then 10,000 cycles using
compression, signed shear, and liner-relevant bending. & 1,000-cycle soft
prosthetic-sensor and $>10{,}000$-cycle printed-PDMS conductor precedents
\cite{q3kwak2020,albrecht2018stretch}. & No open/short, delamination, exposed
silver, or leakage; report $\Delta R/R_0$, calibration shift, hysteresis, and
24~h recovery. Do not substitute continuity for sensor accuracy. \\
\rowcolor{AUBPale}
ENV-01 & Repeat calibration after 35--37~$^{\circ}$C conditioning, 95\%RH,
artificial sweat, and combined cyclic loading. & Prosthetic microclimate and
biofluid exposure are established integration concerns
\cite{paterno2020,q3kwak2020}. & No damage and no $>10\%$ full-scale baseline
shift after defined recovery; failing modalities are reported separately. \\
SAFE-01 & No conductive silver, sharp transition, or fluid path shall contact
skin. Skin-side materials, cleaning method, and adhesives shall be identified
before participant use. & Silicone encapsulation and protected sensor
integration are published precedents \cite{q3kwak2020,paterno2020}; this is an
engineering prototype, not yet cleared for human use. & Dry/wet insulation,
visual/tactile inspection, documented material list, and institutional safety
review before any participant protocol. \\
\rowcolor{AUBPale}
MFG-02 & Every slab shall carry a traveler recording material batch, mix ratio,
degassing, cure, plasma settings, print delay, ink batch, print passes, drying,
encapsulation, and inspection. & Direct-PDMS studies show sensitivity to
surface treatment, geometry, precure, and solvent interaction
\cite{albrecht2018stretch,sun2016,q4abukhalaf2019}. & No missing process fields;
specimen ID links traveler, calibration files, and CAD coordinate manifest. \\
MFG-03 & A bill of materials shall list quantity, supplier, unit cost, reusable
equipment, AUB availability, and fabrication time; cost is reported, not
guessed. & \textbf{D/V:} directly answers Q1 manufacturability; literature
cannot establish local inventory or 2026 pricing. & Completed department
availability table, receipts/quotes, labour log, yield, and cost per passing
slab. \\
\rowcolor{AUBPale}
SYS-01 & A slab passes only if geometry, electrical continuity, modality
calibration, cross-sensitivity, coordinate registration, environmental
recovery, and data integrity all pass. & \textbf{D:} system-level gate derived
from Q1 and the eight Q3 functional blocks. & Signed requirements matrix with
P/D/V status, raw evidence path, pass/partial/fail result, and corrective
action. \\
\end{tabularx}
\endgroup
\sourceanchor{Durability precedents from printed PDMS conductors and soft
prosthetic-interface sensors \cite{albrecht2018stretch,q3kwak2020}; process
traceability is required because those source processes are not identical to
the present build.}

\vspace{0.35em}
\begingroup
\setlength{\fboxsep}{7pt}
\fcolorbox{AUBRed}{AUBPale}{%
\begin{minipage}{0.96\linewidth}
{\fontsize{11.6}{13.6}\selectfont\bfseries\color{AUBRed}
V = REQUIREMENTS DEFINED; COMPLETE 5.2 mm SYSTEM NOT YET VERIFIED}
\par\vspace{0.28em}
{\fontsize{8.6}{10.1}\selectfont
The literature supports every major functional mechanism and supplies concrete
starting processes, geometries, ranges, and validation methods. It does
\textbf{not} supply a completed equivalent of the present slab. The final Q4
decision is therefore to fabricate the 5.2~mm H0 PDMS slab under a recorded
process window, apply the modality and system gates above, and advance to
curved phantom or whole-liner integration only after the flat specimen passes.
Hydrogel remains a later controlled comparison; the map remains an
\emph{interface-condition map} until participant evidence supports a discomfort
or clinical-risk interpretation.}
\end{minipage}}
\endgroup

\clearpage
